% Section 6 -- Computational Bridge
% Target: 1 page.

DLC's symbolic non-splicing property is verified twice
(Section~\ref{sec:symbolic}); the computational reduction is the
load-bearing link to real cryptographic assumptions.

\subsection{The EUF-CMA reduction}

\begin{theorem}[T2 -- computational direction, sketch]
Let $A$ be a PPT adversary against DLC's symbolic non-forgery game
(\texttt{SymbolicSayForgery} in \texttt{models/easycrypt/Game.eca}).
There exists a PPT reduction adversary $A^*$ against the standard
EUF-CMA game (\texttt{EUF\_CMA}) such that:
\[
\Pr[\text{SymbolicSayForgery}(A) = 1] \leq \Pr[\text{EUF\_CMA}(A^*) = 1]
\]
By the EUF-CMA security of Ed25519, the right-hand side is
negligible, hence so is the left.
\end{theorem}

\paragraph{Game-hop sequence (sketched).}
\begin{enumerate}
\item $G_0$: \texttt{SymbolicSayForgery}.  Adversary $A$ given access to
  a signing oracle for a fresh keypair must produce a $(m^*, \sigma^*)$
  pair with $m^*$ not in the queried set and $\sigma^*$ verifying.
\item $G_1$: lazy-sample the signing oracle into a table.  Standard
  identical-until-bad rewrite.
\item $G_2$: replace the signing key with an EUF-CMA challenger.
  The reduction $A^*$ simulates $A$'s oracle by forwarding to its own
  EUF-CMA oracle.  $A$'s eventual forgery is also an EUF-CMA
  forgery.
\item Conclude: $\Pr[G_0(A)] = \Pr[G_2(A^*)] \leq \mathsf{Adv}^{\mathsf{EUF-CMA}}_{\mathsf{Ed25519}}(A^*)$.
\end{enumerate}

\paragraph{Mechanization path.}
Two viable routes:

\begin{enumerate}
\item \textbf{EasyCrypt.}  The skeleton in
  \texttt{models/easycrypt/Game.eca} is structurally complete: type
  declarations, signing oracles, the EUF-CMA module, the
  \texttt{SymbolicSayForgery} module, the EUF-CMA assumption as an
  axiom, and the reduction lemma documented (proof body absent;
  CI parse-checks via \texttt{easycrypt -batch}).  Closure
  requires cryptographer pairing (Blanchet was the named
  collaborator in our plan).

\item \textbf{VCVio.}  Tuma et al.~(2026)~\cite{tuma2026vcvio}
  mechanize EUF-CMA for Schnorr signatures in Lean 4 using a
  monadic oracle-effects framework.  Their proof template
  ports directly to Ed25519 (or to an abstract EUF-CMA-secure
  signature scheme).  This route keeps T2 inside the Lean
  development that holds T1, T3, and T4.
\end{enumerate}

Both routes deliver the same theorem.  VCVio is the autonomously
preferred route as of May 2026 because it stays in one proof
assistant; EasyCrypt is the publication-standard route because of
its cryptographer community.

\paragraph{Status.} The EasyCrypt skeleton parses against the
current grammar (CI step \texttt{easycrypt -batch Game.eca} passes,
6m36s).  The full reduction proof is open work.
